\section*{Chapter \ref{ch:b1:titration-methods} --- Titration Methods and Curves}

\begin{solution}{exo:b1:titration-methods:1}
\ce{H3O+ + OH- -> 2H2O}, $K = 1/K_e = 10^{14.00}$. $C = 12.4 \times
0.100/20.0 = \qty{0.0620}{mol/L}$.
\end{solution}

\begin{solution}{exo:b1:titration-methods:2}
At \omterm{def:b1:titration-methods:titration}{equivalence} the solution is ammonium chloride at
\qty{0.050}{mol/L}: $\mathrm{pH} = \frac12(9.25 - \log 0.050) = 5.28$.
Methyl red (3.8--5.8) contains it; methyl orange, whose zone (2.5--4.5) is reached only after \omterm{def:b1:titration-methods:titration}{equivalence} as the pH falls, would change too late.
\end{solution}

\begin{solution}{exo:b1:titration-methods:3}
\ce{MnO4- + 5Fe^2+ + 8H+ -> Mn^2+ + 5Fe^3+ + 4H2O};
$C = 5 \times 0.0200 \times 12.6/10.0 = \qty{0.126}{mol/L}$.
\end{solution}

\begin{solution}{exo:b1:titration-methods:4}
$C = 0.0100 \times 14.2/50.0 = \qty{2.84e-3}{mol/L}$ of calcium plus
magnesium. Both react before the \omterm{def:b1:titration-methods:titration}{end point}, which gives only their sum:
a \omterm{def:b1:titration-methods:kinds}{simultaneous titration}.
\end{solution}

\begin{solution}{exo:b1:titration-methods:5}
0~mL: \omterm{def:b1:predominant-reaction:strong-acid}{weak acid}, $\mathrm{pH} = \frac12(4.76 + 1.00) = 2.88$. 5.0~mL:
\omterm{def:b1:titration-methods:titration}{half-equivalence}, 4.76. 10.0~mL: 8.73
(\cref{prop:b1:titration-methods:ph-at-equivalence}). 15.0~mL: excess
hydroxide \qty{0.50}{mmol} in \qty{25.0}{mL}, $[\ce{OH-}] = 0.020$,
$\mathrm{pH} = 12.30$. All agree with the computed curve to 0.01.
\end{solution}

\begin{solution}{exo:b1:titration-methods:6}
The $\mathrm{p}K_a$ 2.15 and 7.21 differ by more than 4, as do 7.21 and
12.34: two \omterm{def:b1:titration-methods:titration}{equivalences}, at 10.0 and \qty{20.0}{mL}. First: solution of
\ce{H2PO4-}, $\mathrm{pH} \approx \frac12(2.15 + 7.21) = 4.68$; second:
\ce{HPO4^2-}, $\frac12(7.21 + 12.34) = 9.78$ (the exact curve gives 4.71 and
9.66, the formulas neglecting dilution and water). The third acidity
($\mathrm{p}K_a$ 12.34) is too close to that of water: the hydroxide added
reacts only partly and the pH rises without a jump.
\end{solution}

\begin{solution}{exo:b1:titration-methods:7}
Hydrochloric acid: $6.2 \times 0.100/10.0 = \qty{0.062}{mol/L}$; ethanoic
acid: $(14.6 - 6.2) \times 0.100/10.0 = \qty{0.084}{mol/L}$. Halfway, half
of the ethanoic acid is titrated: $\mathrm{pH} = \mathrm{p}K_a = 4.76$.
\end{solution}

\begin{solution}{exo:b1:titration-methods:8}
$V_{\mathrm{eq}} = \qty{10.0}{mL}$. 2.0~mL: $E = 0.77 + 0.059\log(2/8) =
0.73$~V; 9.0~mL: $0.77 + 0.059\log 9 = 0.83$~V; 11.0~mL: excess
permanganate \qty{0.020}{mmol}, \ce{Mn^2+} \qty{0.200}{mmol}, $E = 1.51 +
\frac{0.059}{5}\log 0.10 = 1.50$~V; 20.0~mL: 1.51~V.
\end{solution}

\begin{solution}{exo:b1:titration-methods:9}
0~mL: $\kappa = (349.8 + 76.2) \times 0.0100 = \qty{4.26}{mS/cm}$.
10.0~mL: \ce{Na+} and \ce{Cl-} at $1.00/110 = \qty{9.09e-3}{mol/L}$,
$\kappa = (50.3 + 76.2) \times \num{9.09e-3} = 1.15$. 20.0~mL: \ce{Na+}
$0.0167$, \ce{Cl-} and \ce{OH-} $0.00833$~mol/L, $\kappa = 0.84 + 0.635 + 1.65
= 3.13$~mS/cm. Slopes (dilution neglected, \qty{1}{mL} adds
\qty{1.0e-3}{mol/L}): $(50.3 - 349.8) \times 10^{-3} = -0.30$ and
$(50.3 + 198.3) \times 10^{-3} = +0.25$~mS/cm per mL.
\end{solution}

\begin{solution}{exo:b1:titration-methods:10}
At $V_{\mathrm{eq}} - v$ the excess acid is $C_Bv$, at $V_{\mathrm{eq}} + v$
the excess base is $C_Bv$; in the same volume, $[\ce{H3O+}]$ before equals
$[\ce{OH-}]$ after, so $\mathrm{pH}(V_{\mathrm{eq}} - v) +
\mathrm{pH}(V_{\mathrm{eq}} + v) = 14$: symmetry about (V$_{\mathrm{eq}}$,
7). Jump: \qty{1.0e-5}{mol} excess in \qty{20}{mL}, pH 3.30 to 10.70, 7.4
units. Divided by 100: $\num{1.0e-7}$~mol in \qty{20}{mL}, pH 5.3 to 8.7,
3.4 units: a much less sharp \omterm{def:b1:titration-methods:titration}{end point}.
\end{solution}

\begin{solution}{exo:b1:titration-methods:11}
Before: \ce{Fe^3+} formed $= 5C_{\mathrm{Mn}}V$, \ce{Fe^2+} left $=
5C_{\mathrm{Mn}}(V_{\mathrm{eq}} - V)$, so $E = 0.77 + 0.059\log\frac{V}{V_{\mathrm{eq}}
- V}$: at $V_{\mathrm{eq}}/2$, $E = 0.77$. After: excess \ce{MnO4-} $\propto
V - V_{\mathrm{eq}}$, \ce{Mn^2+} $\propto V_{\mathrm{eq}}$: $E = 1.51 +
\frac{0.059}{5}\log\frac{V - V_{\mathrm{eq}}}{V_{\mathrm{eq}}}$ (pH 0); at
$2V_{\mathrm{eq}}$, 1.51~V. At \omterm{def:b1:titration-methods:titration}{equivalence} $(0.77 + 5 \times 1.51)/6 =
1.39$~V. At pH 1 the permanganate couple is lower by $0.059 \times 8/5 =
0.094$~V: $E_{\mathrm{eq}} = (0.77 + 5 \times 1.416)/6 = 1.31$~V.
\end{solution}

\begin{solution}{exo:b1:titration-methods:12}
\ce{NH4+ + OH- -> NH3 + H2O} has $K = 10^{14.00 - 9.25} = 10^{4.75}$: the
reaction is not complete near the \omterm{def:b1:titration-methods:titration}{equivalence} and the jump is small,
around pH 11 (\omterm{def:b1:titration-methods:titration}{equivalence} pH $14 - \frac12(4.75 + 1.30) = 11.0$). \omterm{def:b1:titration-methods:kinds}{Back
titration}: add $n_1$ of sodium hydroxide (excess), boil to drive off the
ammonia formed, then titrate the hydroxide left with hydrochloric acid
($n_2$, sharp strong--strong jump): $n(\ce{NH4+}) = n_1 - n_2$.
\end{solution}

\begin{solution}{pb:b1:titration-methods:1}
\textbf{1.} \ce{C6H6O6 + 2H+ + 2e- -> C6H8O6}.
\textbf{2.} \ce{C6H8O6 + I2 -> C6H6O6 + 2I- + 2H+}.
\textbf{3.} 0 and $-$I.
\textbf{4.} \ce{I2 + 2S2O3^2- -> 2I- + S4O6^2-}, $\log K = 2(0.53 - 0.02)/0.059
= 17$.
\textbf{5.} The blue colour needs iodine; it disappears with the last
iodine, at the \omterm{def:b1:titration-methods:titration}{equivalence}.
\textbf{6.} Losses to air would lower the result; added at once and in
excess, the iodine oxidises the ascorbic acid quickly and completely.
\textbf{7.} A fast, \omterm{def:b1:extent-q-and-k:quantitative}{quantitative reaction} during the whole \omterm{def:b1:titration-methods:titration}{titration},
without air oxidation while the \omterm{def:b1:titration-methods:titration}{titrant} is slowly added, and an \omterm{def:b1:titration-methods:titration}{end point}
at the first excess of iodine.
\textbf{8.} Excess iodine $n_R$; reaction with ascorbic acid; \omterm{def:b1:titration-methods:titration}{titration} of
the iodine left by thiosulfate.
\textbf{9.} Otherwise some ascorbic acid would remain and the iodine
excess, zero, would not tell how much. Here
\qty{2.775e-4}{mol} reacted out of \qty{5.000e-4}{mol}: excess indeed.
\textbf{10.} Thiosulfate solutions change slowly with time; its
concentration enters the result directly.
\textbf{11.} 10.
\textbf{12.} A \qty{20.00}{mL} volumetric pipette.
\textbf{13.} $20.00 \times 10^{-3} \times 0.0250 = \qty{5.000e-4}{mol}$.
\textbf{14.} $8.90 \times 10^{-3} \times 0.0500 = \qty{4.450e-4}{mol}$.
\textbf{15.} Half: \qty{2.225e-4}{mol}.
\textbf{16.} $5.000 - 2.225 = \qty{2.775e-4}{mol}$.
\textbf{17.} One to one: \qty{2.775e-4}{mol}.
\textbf{18.} $\times 10$: \qty{2.775e-3}{mol}, $\times 176.0$:
\qty{0.488}{g}.
\textbf{19.} About \qty{2}{\percent} below the label.
\textbf{20.} $n = \bigl(C_IV_I - \frac12C_TV_T\bigr)\,V_{\mathrm{flask}}/V_{\mathrm{aliquot}}$.
\textbf{21.} Iodine introduced: pipette $0.03/20.00 = \qty{0.15}{\percent}$
and concentration \qty{0.2}{\percent}, about \qty{0.25}{\percent}, that is
\qty{1.3e-6}{mol}. Excess: burette $0.05/8.90 = \qty{0.56}{\percent}$ and
concentration \qty{0.2}{\percent}, about \qty{0.60}{\percent}, that is
\qty{1.3e-6}{mol}.
\textbf{22.} The absolute uncertainties combine, while the difference is
smaller than either amount: $\sqrt{1.3^2 + 1.3^2} \times 10^{-6} =
\qty{1.8e-6}{mol}$, \qty{0.66}{\percent} of \qty{2.775e-4}{mol}, more than
either term.
\textbf{23.} With the flask (\qty{0.1}{\percent}) and the aliquot pipette
(\qty{0.2}{\percent}): $\sqrt{0.66^2 + 0.1^2 + 0.2^2} = \qty{0.70}{\percent}$.
\textbf{24.} $0.488 \pm 0.003$~g: \textbf{\qty{0.49}{g} of ascorbic acid
per tablet}.
\end{solution}
